\documentclass[10pt,a4paper]{amsart}
\usepackage[margin=19mm]{geometry}
\usepackage{amsmath,amssymb,mathtools}
\usepackage[scr=rsfs,cal=euler]{mathalfa}
\usepackage{xfrac}
\usepackage[unicode,hidelinks,bookmarks=false]{hyperref}
\newcommand{\ie}{i.e.\ }
\newcommand{\cf}{cf.\ }
\newcommand{\resp}{respectively\ }
\newcommand{\Subsection}{\subsection}
\providecommand{\nobreakdash}{\nobreak\hbox{-}}
\setlength{\emergencystretch}{2em}
\multlinegap0pt
\usepackage{bm}
\usepackage{bbm}
\usepackage{dsfont}
\usepackage{xspace}
\usepackage{cancel}
\usepackage{mathtools}
\usepackage[shortlabels]{enumitem}
\usepackage{amsthm}
\makeatletter
\@ifundefined{assumption}{\newtheorem{assumption}{Assumption}}{}
\@ifundefined{definition}{\newtheorem{definition}{Definition}}{}
\@ifundefined{lemma}{\newtheorem{lemma}{Lemma}}{}
\@ifundefined{theorem}{\newtheorem{theorem}{Theorem}}{}
\makeatother
\newcommand{\Acal}{\mathcal{A}}
\newcommand{\Bcal}{\mathcal{B}}
\newcommand{\EE}{\mathbb{E}}
\newcommand{\PP}{\mathbb{P}}
\newcommand{\RR}{\mathbb{R}}
\newcommand{\Scal}{\mathcal{S}}
\newcommand{\afrak}{\mathfrak{a}}
\newcommand{\barf}{\overline{f}}
\newcommand{\hC}{\hat{C}}
\newcommand{\hD}{\hat{D}}
\newcommand{\xstar}{x^*}
\newcommand{\ystar}{y^*}
\title[Finite-Time Bounds for Two-Time-Scale Stochastic Approximati]{Finite-Time Bounds for Two-Time-Scale Stochastic Approximation with Arbitrary Norm Contractions and Markovian Noise}
\author{Siddharth Chandak, Shaan Ul Haque, Nicholas Bambos}
\date{Source: arXiv:2503.18391v2 (CC BY 4.0)}
\begin{document}
\maketitle

\noindent\emph{Source and licence.} Finite-Time Bounds for Two-Time-Scale Stochastic Approximation with Arbitrary Norm Contractions and Markovian Noise. Version arXiv:2503.18391v2 (CC BY 4.0), distributed under the Creative Commons Attribution 4.0 International licence, as recorded in the arXiv licence metadata for this exact version and shown on the abstract page https://arxiv.org/abs/2503.18391v2. Full rights notice: licenses/arxiv-2503.18391-CC-BY-4.0.txt.\par\smallskip

\noindent\textit{Editorial scope.} Counted unit: the Lemma whose source label is \texttt{lem:inter-bounds} in the article (source0.tex lines 1319--1336, source label \texttt{lem:inter-bounds}), delivered together with the article's own complete original proof of it (source0.tex lines 1337--1380, one \texttt{proof} environment of 44 lines). No mathematics is written, repaired or re-derived: statement and proof are the article's own text line for line. Required context retained uncounted: assu-Markov (lines 208--215); assu-contractive-f (lines 221--227); assu-Lipschitz (lines 243--252); assu-stepsize (lines 255--258); thm:expectation-bound (lines 266--272); defn:Moreau (lines 303--310); defn:VandW (lines 328--335); lem:xstar-Lipschitz (lines 354--360); lemma:subg (lines 780--788); lem:poisson\_soln-prop (lines 799--806). Every definition, display and label of those lines is retained unchanged. Omitted from this package and not redistributed: 1--207, 216--220, 228--242, 253--254, 259--265, 273--302, 311--327, 336--353, 361--779, 789--798, 807--1318, 1381--1676. Nothing inside a retained excerpt is omitted or altered: every retained body is delivered exactly as the article's lines are written, so no omission receipt of any kind accompanies this package. Citations: the retained excerpts carry no citation command at all, so no bibliography entry of the article is transcribed and no reference is redistributed. Reference adaptation: none was needed and none was made; every reference inside the delivered unit resolves inside it, so no unresolved reference or citation is produced. Printed slips kept literal: no printed word, symbol, hypothesis or number of the counted statement or of its proof was altered. Layout and class: the article's own class is \texttt{ieeeconf} with packages including amsmath, amssymb, color, cite, graphicx, latexsym, algorithm, algpseudocode, verbatim, babel, balance, amsthm, enumitem, xcolor; the wrapper is the lane's pinned \texttt{amsart} editorial wrapper at 10pt on A4, which restates the theorem-like environments the retained text needs and adds the article's own macro definitions that the retained text needs (\texttt{\textbackslash Acal}, \texttt{\textbackslash Bcal}, \texttt{\textbackslash EE}, \texttt{\textbackslash PP}, \texttt{\textbackslash RR}, \texttt{\textbackslash Scal}, \texttt{\textbackslash afrak}, \texttt{\textbackslash barf}, \texttt{\textbackslash hC}, \texttt{\textbackslash hD}, \texttt{\textbackslash xstar}, \texttt{\textbackslash ystar}), with extra wrapper packages (bm, bbm, dsfont, xspace, cancel, mathtools, enumitem). The delivered numbering is the unit's own. Assets: no retained span contains a figure environment, a graphics-inclusion command, a PDF-page inclusion, a table or a TikZ picture; no image file is copied into this package, the package declares no asset dependency and no image file is redistributed with it. Licence: version arXiv:2503.18391v2 of this article is distributed under the Creative Commons Attribution 4.0 International licence, as recorded in the arXiv licence metadata for this exact version and shown on the abstract page https://arxiv.org/abs/2503.18391v2. A full rights notice is delivered at licenses/arxiv-2503.18391-CC-BY-4.0.txt. Characters: every retained excerpt is pure ASCII, so the delivered engine, XeLaTeX with its default Latin Modern text font, reports no missing character.
\begin{assumption}\label{assu-Markov}
    The random process $\{Z_n\}$ is a Markov chain taking values in a finite state space $\Scal$ and satisfies
    \begin{align*}
        \PP(Z_{n+1}=j\mid Z_m, m\leq n)&=\PP(Z_{n+1}=j\mid Z_n)\\
        &=p(j\mid Z_n),
    \end{align*}
    for $j\in\Scal$. $p(\cdot|\cdot)$ is the transition probability of an irreducible Markov chain with unique stationary distribution $\pi(\cdot)$.
\end{assumption}

\begin{assumption}\label{assu-contractive-f}
Function $\barf(x,y)$ is $\lambda$-contractive in $x$, i.e., there exists a $\lambda\in [0,1)$ such that 
\begin{equation}
    \|\barf(x_1,y)-\barf(x_2,y)\|\leq \lambda \|x_1-x_2\|,
\end{equation}
for all $x_1,x_2\in\RR^{d_1}, y\in\RR^{d_2}$.
\end{assumption}

\begin{assumption}\label{assu-Lipschitz}
Functions $f(\cdot)$ and $g(\cdot)$ are $L$-Lipschitz, i.e.,
\begin{align*}
    &\|f(x_1,y_1,i)-f(x_2,y_2,i)\|\leq L(\|x_1-x_2\|+\|y_1-y_2\|)\\
    &\|g(x_1,y_1,i)-g(x_2,y_2,i)\|\leq L(\|x_1-x_2\|+\|y_1-y_2\|),
\end{align*}
for all $x_1,x_2\in\RR^{d_1}, y_1,y_2\in\RR^{d_2}$ and $i\in\Scal$. Moreover,
$$\|f(x,y,i)\|+\|g(x,y,i)\|\leq K(1+\|x\|+\|y\|),$$
for $x\in\RR^{d_1}, y\in\RR^{d_2},i\in\Scal$,
\end{assumption}

\begin{assumption}\label{assu-stepsize}
The step size sequences are of the form $$\alpha_n=\frac{\alpha_0}{(n+1)^{\afrak}}\;\text{and}\;\beta_n=\frac{\beta_0}{n+1},$$
where $0<\alpha_0,\beta_0<1$ and $0.5<\afrak<1$. 
\end{assumption}

\begin{theorem}\label{thm:expectation-bound}
    Suppose that Assumptions \ref{assu-Markov}-\ref{assu-stepsize} hold. Then the iterates $(x_n,y_n)$ almost surely converge to the fixed points $(\xstar,\ystar)$. There exists constant $C_1,C_2,C_3>0$ and time $n_0\geq 0$ such that for all $n\geq n_0$ and $\beta_0\geq C_3$,
    \begin{align*}
        &\EE\left[\|x_n-\xstar\|^2\right]\leq C_1\left(\alpha_n+\frac{\beta_n^2}{\alpha_n^2}\right),\\
    \text{and}\;\;&\EE\left[\|y_n-\ystar\|^2\right]\leq C_2\left(\alpha_n+\frac{\beta_n^2}{\alpha_n^2}\right).
    \end{align*}
\end{theorem}

\begin{definition}\label{defn:Moreau}
    For $x\in\RR^{d_1}$, $y\in\RR^{d_2}$ and $q>0$ define functions $\Acal:\RR^{d_1}\mapsto\RR$ and $\Bcal:\RR^{d_2}\mapsto\RR$ as 
\begin{align*}
\Acal(x)&=\inf_{v\in\RR^{d_1}}\left\{\frac{1}{2}\|v\|^2+\frac{1}{2q}\|x-v\|_2^2\right\}    \\
\Bcal(y)&=\inf_{w\in\RR^{d_2}}\left\{\frac{1}{2}\|w\|^2+\frac{1}{2q}\|y-w\|_2^2\right\}.   
\end{align*}
Here $\|\cdot\|_2$ denotes the $\ell_2$ norm. 
\end{definition}

\begin{definition}\label{defn:VandW}
    For a fixed $i_0\in\Scal$, $\tau\coloneqq \min\{n>0: \!Z_n=i_0\}$ and $E_i[\cdot]=E[\cdot|Z_0=i]$, define $V(x,y,i)$ and $W(x,y,i)$ as
    \begin{align*}
        V(x,y,i)&=E_i\left[\sum_{m=0}^{\tau-1}\Big(f(x,y,Z_m)-\sum_{j\in\Scal}\pi(j)f(x,y,j)\Big)\right],\\
        W(x,y,i)&=E_i\left[\sum_{m=0}^{\tau-1}\Big(g(x,y,Z_m)-\sum_{j\in\Scal}\pi(j)g(x,y,j)\Big)\right],
    \end{align*}
for all $x\in\RR^{d_1},y\in\RR^{d_2}$ and $i\in\Scal$.
\end{definition}

\begin{lemma}\label{lem:xstar-Lipschitz}
    Suppose Assumptions \ref{assu-contractive-f} and \ref{assu-Lipschitz} hold. Then the map $y\mapsto\xstar(y)$ is Lipschitz with parameter $L_1\coloneqq L/(1-\lambda)$, i.e.,
    \begin{equation*}
      \|\xstar(y_1)-\xstar(y_2)\|\leq L_1\|y_1-y_2\|,
    \end{equation*}
    for $y_1,y_2\in\RR^{d_2}$.
\end{lemma}

\begin{lemma}\label{lemma:subg}
    For $\Acal(\cdot)$ and $\Bcal(\cdot)$ as defined in Definition \ref{defn:Moreau}, the following hold true.
    \begin{enumerate}[label=(\alph*)]
        \item  For $x_1,x_2\in\RR^{d_1}$,  $\langle \nabla \Acal(x_1),x_2\rangle_2 \leq \|x_1\|_\Acal\|x_2\|_\Acal.$ Similarly, for $y_1,y_2\in\RR^{d_2}$,  $\langle \nabla \Bcal(y_1),y_2\rangle_2 \leq \|y_1\|_\Bcal\|y_2\|_\Bcal.$
        \item For $x\in\RR^{d_1}$, 
        $\langle \nabla \Acal(x), x\rangle_2\geq 2\Acal(x).$ Similarly, for $y\in\RR^{d_2}$, 
        $\langle \nabla \Bcal(y), y\rangle_2\geq 2\Bcal(y).$
    \end{enumerate}
\end{lemma}

\begin{lemma}\label{lem:poisson_soln-prop}
    The solutions of Poisson equation $V(\cdot,\cdot,\cdot)$ and $W(\cdot,\cdot,\cdot)$ as defined in Definition \ref{defn:VandW} satisfy linear growth and Lipschitzness, i.e., the following hold.
    \begin{enumerate}[label=(\alph*)]
        \item $\|V(x,y,i)\|+\|W(x,y,i)\|\leq K'(1+\|x\|+\|y\|)\; \forall x\in\RR^{d_1}, y\in\RR^{d_2}, i\in\Scal.$
        \item $\|V(x_1,y_1,i)-V(x_2,y_2,i)\|\leq L_2(\|x_1-x_2\|+\|y_1-y_2\|) \; \forall x_1,x_2\in\RR^{d_1}, y_1,y_2\in\RR^{d_2}, i\in\Scal.$
        \item $\|W(x_1,y_1,i)-W(x_2,y_2,i)\|\leq L_2(\|x_1-x_2\|+\|y_1-y_2\|) \; \forall x_1,x_2\in\RR^{d_1}, y_1,y_2\in\RR^{d_2}, i\in\Scal.$
    \end{enumerate}
\end{lemma}

\begin{lemma}\label{lem:inter-bounds}
    Suppose the setting of Theorem \ref{thm:expectation-bound} holds. Then, the following hold true.
    \begin{enumerate}[label=(\alph*)]
        \item \label{lem-part:gen-diff-bound} For constants $\hC_1,\hC_2>0$, 
        \begin{align*}
        \hC_1\|x_{n+1}-x_n\|+\hC_2\|y_{n+1}-y_n\|\leq \alpha_n(1+2K)\left(\hC_1+\hC_2\gamma_n\right)(1+\|x_n\|+\|y_n\|).
        \end{align*}
        \item \label{lem-part:x-diff-bound} Define constant $\hD_1=(1+2K)(1+L_1\gamma_0)$. Then,
        \begin{align*}
        \|x_{n+1}-x_n+\xstar(y_n)-\xstar(y_{n+1})\|\leq \alpha_n\hD_1(1+\|x_n\|+\|y_n\|).
        \end{align*} 
        \item \label{lem-part:crude-bound} Define $\hD_2=3\max\{(1+\|x^*\|+\|y^*\|)^2,(1+L_1)^2(1+q/\ell^2)\}$. Then
    \begin{align*}
        (1+\|x_n\|+\|y_n\|)^2\leq \hD_2(1+\Acal(x_n-x^*(y_n))+\Bcal(y_n-y^*)).
    \end{align*}
    \item \label{lem-part:telescoping-diff-bound}Define $\hD_3\coloneqq\frac{K'(2+\hD_2)}{2\sqrt{1+q/\ell^2}}$. Then $$|d_n|\leq \hD_3(1+\Acal(x_n-\xstar(y_n))+\Bcal(y_n-\ystar))\;\text{and}\;|e_n|\leq \hD_3(1+\Acal(x_n-\xstar(y_n))+\Bcal(y_n-\ystar)).$$
    \end{enumerate}
\end{lemma}

\begin{proof}
    \begin{enumerate}[label=(\alph*)]
        \item We have
    \begin{align*}
        \hC_1\|x_{n+1}-x_n\|+\hC_2\|y_{n+1}-y_n\|
        &\leq \Bigg(\alpha_n\hC_1\|f(x_n, y_n, Z_n)-x_n+M_{n+1}\|\\
        &\;\;\;\;\;\;\;\;\;+\beta_n\hC_2\|g(x_n,y_n,Z_n)- y_n+ M'_{n+1}\|\Bigg).
    \end{align*}
    Using the linear growth assumptions
    \begin{align*}
        \hC_1\|x_{n+1}-x_n\|+\hC_2\|y_{n+1}-y_n\|
        &\leq \alpha_n\hC_1((2K)(1+\|x_n\|+\|y_n\|)+\|x_n\|)\\
        &\;\;\;+\beta_n\hC_2(2K)(1+\|x_n\|+y_n\|)+\|y_n\|)\\
        &\leq \alpha_n\left(\hC_1(1+2K)+\hC_2\gamma_n(1+2K)\right)(1+\|x_n\|+\|y_n\|).
    \end{align*}
    \item     Using Lemma \ref{lem:xstar-Lipschitz}, we have
    \begin{align*}
        \|x_{n+1}-x_n+\xstar(y_n)-\xstar(y_{n+1})\|
        &\leq \|x_{n+1}-x_n\|+L_1\|y_n-y_{n+1}\|.
    \end{align*}
    Put $\hC_1=1$ and $\hC_2=L_1$ in the Lemma \ref{lem:inter-bounds}\ref{lem-part:gen-diff-bound} and use the fact that $\gamma_n$ is a non-increasing sequence. The claim follows.
    \item Note that
    \begin{align*}
        (1+\|x_n\|+\|y_n\|)
        &\leq (1+\|x^*(y_n)\|+\|y^*\|+\|x_n-x^*(y_n)\|+\|y_n-y^*\|).
    \end{align*}
    Using Lemma \ref{lem:xstar-Lipschitz}, we can bound $\|x^*(y_n)\|\leq L_1\|y_n-y^*\|+\|x^*\|$ to get
    \begin{align*}
        (1+\|x_n\|+\|y_n\|)^2&\leq 3((1+\|x^*\|+\|y^*\|)^2+\|x_n-x^*(y_n)\|^2+(1+L_1)^2\|y_n-y^*\|^2)\\
        &\leq 3((1+\|x^*\|+\|y^*\|)^2+(1+q/\ell^2)\Acal(x_n-x^*(y_n))\\
        &\;\;\;+(1+L_1)^2(1+q/\ell^2)\Bcal(y_n-y^*)).
    \end{align*}
    The claim follows. 
    \item Using Lemma \ref{lemma:subg}, we have:
    \begin{align*}
        |d_n|&\leq \|x_n-\xstar(y_n)\|_\Acal\|V(x_n,y_n,Z_n)\|_\Acal\\
        &\stackrel{(a)}{\leq} \frac{1}{\sqrt{1+q/\ell^2}}\|x_n-\xstar(y_n)\|_\Acal\|V(x_n,y_n,Z_n)\|\\
        &\stackrel{(b)}{\leq} \frac{K'}{\sqrt{1+q/\ell^2}}\|x_n-\xstar(y_n)\|_\Acal(1+\|x_n\|+\|y_n\|)\\
        &\stackrel{(c)}{\leq} \frac{K'}{\sqrt{1+q/\ell^2}}\left(\frac{1}{2}\|x_n-\xstar(y_n)\|_\Acal^2+\frac{1}{2}(1+\|x_n\|+\|y_n\|)^2\right)\\
        &\stackrel{(d)}{\leq} \frac{K'}{\sqrt{1+q/\ell^2}}\left(\Acal(x_n-\xstar(y_n))+\frac{1}{2}\hD_2\left(1+\Acal(x_n-x^*(y_n)+\Bcal(y_n-y^*)\right)\right)
    \end{align*}
    where we use equivalence of norms for $(a)$, Lemma \ref{lem:poisson_soln-prop} for $(b)$, Young's inequality for $(c)$ and Lemma \ref{lem:inter-bounds}\ref{lem-part:crude-bound} for $(d)$. The bound for $e_n$ follows in exactly the same manner.
    \end{enumerate}
\end{proof}

\end{document}
